% ============================================================================
% Fichier  : partie2/P02_C05_chaine_possession.tex : FORN (LPI342)
% Sources : ISO/IEC 27037 ; ACPO ; RFC 3227. Opération TSANAGA.
% ============================================================================

\chapter{La Chaîne de Possession}
\label{chap:custody}
\chaptermark{La chaîne de possession}

\epigraph{%
    « Une preuve dont on ignore où elle était hier\\
    ne prouve plus rien aujourd'hui. »%
}{Principe du cours}

\niveauChapitre{Fondamental}{Chapitre~\ref{chap:droit}.}

\begin{boxobjectifs}
\begin{itemize}
    \item Définir la chaîne de possession (ou de custody).
    \item Tenir la continuité de la détention d'une pièce.
    \item Sceller, étiqueter, consigner les transferts.
    \item Relier la chaîne de possession à l'opposabilité.
    \item Rédiger une fiche de scellé.
\end{itemize}
\end{boxobjectifs}

\bigskip

% ============================================================================
\section*{Question directrice : comment prouver que le disque analysé est bien celui qui a été saisi ?}
% ============================================================================

Entre la saisie chez M.~NDONGO et l'analyse, le support passe de main en
main. Si l'on ne peut retracer ce parcours, la défense contestera que la
pièce analysée est bien celle saisie. Ce chapitre pose la chaîne de
possession.

\medskip
\textbf{Plan du chapitre.}
\begin{enumerate}
    \item Définir la chaîne (section~\ref{sec:c05-definir}).
    \item Sceller et étiqueter (section~\ref{sec:c05-sceller}).
    \item Consigner les transferts (section~\ref{sec:c05-transferts}).
    \item Chaîne et opposabilité (section~\ref{sec:c05-opposabilite}).
\end{enumerate}

\separateur

% ============================================================================
\section{Définir la Chaîne}
\label{sec:c05-definir}
% ============================================================================

\begin{boxdef}
\termecle{Chaîne de possession} (ou chaîne de custody)~: continuité
documentée de la détention d'une pièce, de la saisie à la présentation,
établissant à tout instant qui la détenait, où, et pour quoi. Toute rupture
(un instant sans détenteur consigné) fragilise la pièce.
\end{boxdef}

% ============================================================================
\section{Sceller et Étiqueter}
\label{sec:c05-sceller}
% ============================================================================

\begin{boxprocedure}
\textbf{Sceller une pièce}
\begin{enumerate}
    \item Identifier la pièce~: numéro de scellé, description, photographie.
    \item Placer sous scellé inviolable, daté et signé.
    \item Calculer l'empreinte (condensat) dès que possible.
    \item Étiqueter~: affaire, pièce, date, agent.
    \item Consigner l'acte dans le registre des scellés.
\end{enumerate}
\end{boxprocedure}

% ============================================================================
\section{Consigner les Transferts}
\label{sec:c05-transferts}
% ============================================================================

\begin{boxobjet}{Scellé n° TSA-2026-07 : clé USB}
\textbf{Nature~:} clé USB saisie au domicile de M.~NDONGO le 20~avril 2026.
\textbf{Scellé~:} TSA-2026-07, inviolable, signé. \textbf{Empreinte~:}
condensat calculé à la réception. \textbf{Transferts~:} parquet $\to$ NYONG
SÉCURITÉ (20/04), coffre (20/04), analyse (22/04). \textbf{Observation~:}
continuité assurée.
\end{boxobjet}

% ============================================================================
\section{Chaîne et Opposabilité}
\label{sec:c05-opposabilite}
% ============================================================================

\begin{boxCRO}
\textbf{La chaîne de possession sert l'opposabilité.} Une pièce à la
détention continue et documentée est opposable \cro{0}{+}{+}~; une pièce dont
la détention présente un trou devient contestable \cro{0}{0}{-}. La chaîne de
possession est la condition première de la recevabilité.
\end{boxCRO}

\begin{boxscenario}{Opération TSANAGA. Un transfert non consigné}
\textbf{Constatation.} La clé USB est restée une nuit sur le bureau d'un
stagiaire, sans consignation.

\medskip
\textbf{Analyse} (confiance élevée). Ce trou, même sans malveillance,
affaiblit l'opposabilité de la clé. M.~BIKORO le consigne honnêtement et
s'appuie davantage sur l'ordinateur, dont la chaîne est intacte. On ne cache
pas une rupture~: on la déclare.
\end{boxscenario}

\begin{boxmetier}{Analyste forensique et expert judiciaire}
\textbf{La chaîne commence à la saisie.} L'expert exige que la chaîne soit
tenue dès la saisie, pas à la réception~: un support mal scellé sur le
terrain arrive déjà fragilisé au laboratoire.
\end{boxmetier}

\separateur

\begin{boxresume}
\textbf{Ce qu'il faut retenir de ce chapitre~:}
\begin{itemize}
    \item \textbf{La chaîne de possession} documente la détention, de la
          saisie à la présentation.
    \item \textbf{On scelle, étiquette, calcule l'empreinte}.
    \item \textbf{On consigne chaque transfert}.
    \item \textbf{Un trou} rend la pièce contestable.
    \item \textbf{On déclare une rupture}, on ne la cache pas.
\end{itemize}
\end{boxresume}

\begin{boxinfo}
\textbf{Points de contrôle de l'élève.} Je sais~: définir la chaîne de
possession~; sceller une pièce~; consigner un transfert~; relier la chaîne à
l'opposabilité~; réagir à une rupture.
\end{boxinfo}

\bigskip

\begin{boxexo}[title={\ding{45}\quad Exercice~5.1. Définir \niveaubase}]
Qu'est-ce que la chaîne de possession~?
\end{boxexo}

\begin{boxexo}[title={\ding{45}\quad Exercice~5.2. Sceller \niveaubase}]
Citez trois éléments d'une fiche de scellé.
\end{boxexo}

\begin{boxexo}[title={\ding{45}\quad Exercice~5.3. L'empreinte \niveaubase}]
Pourquoi calculer l'empreinte dès la réception~?
\end{boxexo}

\begin{boxexo}[title={\ding{45}\quad Exercice~5.4. Le trou \niveaumoyen}]
La clé USB est restée une nuit sans consignation. Quel effet~?
\end{boxexo}

\begin{boxexo}[title={\ding{45}\quad Exercice~5.5. Déclarer \niveaumoyen}]
Pourquoi déclarer une rupture plutôt que la cacher~?
\end{boxexo}

\begin{boxexo}[title={\ding{45}\quad Exercice~5.6. CRO \niveaumoyen}]
Expliquez l'analyse CRO d'une chaîne continue et d'une chaîne rompue.
\end{boxexo}

\begin{boxexo}[title={\ding{45}\quad Exercice~5.7. Terrain \niveaumoyen}]
Pourquoi la chaîne doit-elle commencer à la saisie, pas à la réception~?
\end{boxexo}

\begin{boxexo}[title={\ding{45}\quad Exercice~5.8. Identité de la pièce \niveaumoyen}]
Comment prouver que le disque analysé est bien celui saisi~?
\end{boxexo}

\begin{boxexo}[title={\ding{45}\quad Exercice~5.9. Fiche de scellé \niveauexpert}]
Rédigez la fiche de scellé complète de la clé USB TSA-2026-07.
\end{boxexo}

\begin{boxexo}[title={\ding{45}\quad Exercice~5.10. Registre \niveauexpert}]
Concevez le registre des scellés du cabinet NYONG SÉCURITÉ.
\end{boxexo}

\bigskip

\begin{boxinfo}
\textbf{Pour aller plus loin.}
\begin{itemize}
    \item ISO/IEC 27037~\cite{iso27037}~; ACPO~\cite{acpo2012}~;
          RFC~3227~\cite{rfc3227}.
\end{itemize}
\end{boxinfo}

\begin{boxinfo}
\textbf{Transition.} Chaîne et droit convergent sur la question de la
recevabilité~: le Chapitre~\ref{chap:recevabilite} la traite.
\end{boxinfo}
